% Auto-generated complete chapter from Research/deep-research-report_00.md
\chapter{Resumen Ejecutivo}
\label{complete:research_053}

\begin{sourcemeta}
\textbf{Fuente:} \href{file:///E:/Laboral/Research/deep-research-report_00.md}{Research/deep-research-report\_00.md}\\
\textbf{Seccion:} research\\
\textbf{Estado:} research\_snapshot\\
\textbf{Rol:} Research / estudio de mercado\\
\textbf{Lineas originales:} 461\\
\textbf{Ultima modificacion:} 2026-06-11 12:18\\
\textbf{Deduplicacion conservadora:} 0 bloques repetidos omitidos; 0 caracteres repetidos no reimpresos.
\end{sourcemeta}

\section*{Resumen de orientacion}
Alexander Woodcock Salomón es un \textbf{Technical Artist y Desarrollador 3D en Tiempo Real} (Unity, WebGL, C\#, Blender, optimización CAD-realidad) con fuerte formación en ingeniería multimedia y electrónica. Su meta es asegurar \textbf{empleos mejor remunerados} (idealmente remotos desde Colombia) en los próximos 3–6 meses, con un plan de 12–24 meses para duplicar o triplicar su salario y lograr movilidad internacional. Dada su experiencia académica (tesis sobre visualización técnica 

\section*{Contenido completo depurado}
\addcontentsline{toc}{section}{Contenido completo depurado}




Alexander Woodcock Salomón es un \textbf{Technical Artist y Desarrollador 3D en Tiempo Real} (Unity, WebGL, C\#, Blender, optimización CAD-realidad) con fuerte formación en ingeniería multimedia y electrónica. Su meta es asegurar \textbf{empleos mejor remunerados} (idealmente remotos desde Colombia) en los próximos 3–6 meses, con un plan de 12–24 meses para duplicar o triplicar su salario y lograr movilidad internacional. Dada su experiencia académica (tesis sobre visualización técnica de dron en Unity WebGL) y sus proyectos con IA asistida, el enfoque de búsqueda se orientará a roles de \emph{Technical Artist, 3D interactivo, XR/simulación/digital twin} y \emph{desarrollo de herramientas Python/IA}, en industrias de simulación industrial, visualización técnica, videojuegos y tecnologías creativas.




Los objetivos salariales de USD 1.5k, 3k y 6k mensuales corresponden a rangos desde niveles locales bajos hasta medianos-globales. Por ejemplo, el salario medio remoto de un \emph{Technical Artist} es de \textasciitilde{}\$90k anuales (\textasciitilde{}\$7.5k/mes). Así, \$1.5k/mes (\$18k/año) queda por debajo de ese promedio global, \$3k/mes (\$36k/año) se acerca a roles junior intermedios internacionales, y \$6k/mes (\$72k/año) es realista para posiciones mid-level en mercados maduros. Estos datos indican que la búsqueda debe priorizar ofertas globales mejor remuneradas (EE.UU./UE) mientras se aprovechan los canales de contratación remota desde Colombia.




En Colombia, el salario típico de un ingeniero de software (\textasciitilde{}COP 123 M anuales, \textasciitilde{}USD 31k) es inferior a los objetivos internacionales. Por tanto, el plan se centra en \textbf{trabajo remoto/contrato B2B internacional}. Se preparará el candidato para contratación remota: obtener RUT, afiliarse a EPS/ARL, usar facturación electrónica DIAN, y manejar pagos internacionales (Wise/Payoneer).




Las barreras legales actuales impiden asumir empleo en la UE sin permiso: ni residencia alemana ni ciudadanía portuguesa están vigentes. Aunque el candidato espera futuro pasaporte portugués por linaje (requiere 3 años de residencia según ley reciente) y casarse en Alemania (facilitaría visa familiar sin idioma), estos son logros a mediano plazo (>2 años). De momento, la prioridad es el trabajo remoto desde Colombia. Se explorarán puestos freelance/contratista mediante EOR o plataformas globales, garantizando cumplimiento fiscal colombiano.




En resumen, la estrategia combina:


\begin{itemize}
\item \textbf{Investigación de mercado y priorización de roles} (Technical Artist, 3D interactivo, XR/simulación, Python/AI) con ajuste al perfil técnico real.  
\item \textbf{Estrategia de CV/LinkedIn/GitHub/portafolio} basada en proyectos clave (p.ej. TwinSight X500) y palabras clave ATS (Unity, C\#, WebGL, optimización 3D, Blender, Python).  
\item \textbf{Plan 30/60/90 días} concreto: limpiar y publicar proyectos relevantes (GitHub, ArtStation), optimizar perfiles profesionales, iniciar contactos y aplicaciones.  
\item \textbf{Plan de acciones de networking y entrevistas}: secuencias de outreach en inglés y español, preparación de respuestas a temas sensibles (baja experiencia formal, uso de IA en proyectos) y estrategia de negociación basada en datos de salario global.  
\end{itemize}




Cada recomendación se fundamenta en datos recientes (2024–2026) y normas locales. Por ejemplo, para empleo remoto desde Colombia se considerarán modelos de contratación B2B/contratista documentando IVA exportación (0\%) y usando pagos internacionales formales. En la ruta internacional, se evitará sobrevalorar las oportunidades europeas inmediatas: la residencia portuguesa exige 3 años en Portugal y la ruta alemana requiere matrimonio y solicitud de visa (sin dominio del alemán). La búsqueda inicial se enfocará en empresas globales remotas de EE.UU., Canadá, Reino Unido, Países Bajos y Australia con alta contratación remota en tech art, XR y simulación.




\section{Diagnóstico de la posición de mercado actual}




El perfil de Alexander combina \textbf{ingeniería multimedia, herramientas 3D y Python/IA}. A nivel laboral:




\begin{itemize}
\item \textbf{Experiencia:} proyectos de portfolio/académicos (p.ej. prototipo WebGL de dron, herramientas de automatización AI). Aún no hay experiencia formal de industria en technical art/Unity, sólo freelance y proyectos personales. Esto es común en candidatos junior de campos emergentes.
\end{itemize}




\begin{itemize}
\item \textbf{Habilidades destacadas:} Unity/C\#, WebGL/URP, Blender (retopología y optimización CAD), UI Toolkit, Shader Graph, Python, LLM y pipeline automation. Estos responden a demandas de industrias de \emph{visualización técnica y simulación}. Sin embargo, debe validarse experiencia con frameworks web (FastAPI, React) u otras tecnologías secundarias si se postula a full-stack o PyDev.
\end{itemize}




\begin{itemize}
\item \textbf{Idiomas:} español nativo, inglés C1, alemán nulo. Esto califica para roles internacionales en inglés; alemán ofrece poca ventaja en corto plazo.
\end{itemize}




\begin{itemize}
\item \textbf{Ubicación:} Colombia – mercado local limitado en tech art (salarios bajos, pocas grandes empresas en videojuegos). El candidato apunta a mercado global remoto: USA, Canadá, Europa, LATAM nearshore. 
\end{itemize}




\begin{itemize}
\item \textbf{Legal:} No tiene permiso UE inmediato. Portugal por linaje podría llegar, pero requiere residir 3 años tras el cambio de ley. La ruta alemana por cónyuge (futura) permite estadía/sitio de trabajo en Alemania sin alemán, pero esto es incierto y a mediano plazo.
\end{itemize}




\textbf{Entorno salarial:} En EEUU/UE, salarios para \emph{Technical Artists} y \emph{Desarrolladores 3D/AR/VR} rondan USD 80k–130k anuales. Es decir, unos USD 6.7k–10.8k mensuales. OBJETIVOS SALARIALES:


\begin{itemize}
\item \textbf{USD 1.5k/mes (\$18k/año):} muy bajo comparado a niveles mundiales; situaría a Alexander por debajo de la mediana local (≈USD 21k–31k/año para devs en Colombia). Quizá factible solo en proyectos freelance locales o pasantías.  
\item \textbf{USD 3k/mes (\$36k/año):} aún bajo para EE.UU./Europa, equivale a nivel muy junior. Sin embargo, está por encima de salarios colombianos medios para devs (≈USD 27k–41k). Probable en el corto plazo como contratista B2B remota en startups o empresas nearshore.  
\item \textbf{USD 6k/mes (\$72k/año):} rango medio-alto global. Muy plausible para mid-level Technical Artist o XR/Sim Developer en mercados clave. Meta final ambiciosa pero alcanzable.
\end{itemize}




\textbf{Tendencias de mercado (2024-2026):} La demanda de perfiles Unity/AR/VR y visualización 3D sigue activa por la industria de videojuegos, metaverso, simulación militar/industrial y sectores de visualización (salud, ingeniería). Empresas remotas están creciendo, y portales de empleo remoto listan numerosas vacantes para \emph{Unity Technical Artist} y \emph{3D Developer}. Por ejemplo, Indeed y ZipRecruiter muestran centenares de ofertas remotas para "Unity Technical Artist" con salarios medios de USD 70k–140k. Sin embargo, la competencia es alta (Glassdoor indica \textasciitilde{}47 candidatos por oferta de Unity dev en 2025). Diferenciarse requerirá portafolio sólido y contactos de la industria.




\textbf{Resumen del diagnóstico:} Alexander tiene un perfil técnico fuerte pero junior, con proyectos relevantes (TwinSight X500) pero sin experiencia corporativa significativa. Su costo local es bajo, su potencial salarial global es alto. La estrategia debe resaltar sus logros académicos/prácticos, posicionarlo como candidato capacitado en Technical Art/3D interactivo, y conectar con el mercado remoto antes de depender de vías migratorias. Se recomienda una búsqueda priorizando empresas globales con mentalidad remota y proyectos de visualización técnica, ajustando expectativas de salario inicial (c. \$18–36k/año), mientras se avanza hacia posiciones mejor pagadas (\$72k/año) en 12-24 meses.




\section{Declaración de posicionamiento recomendada}




Basado en habilidades y proyectos clave, se propone la siguiente declaración de posicionamiento:




\begin{quote}
\textbf{“Technical Artist \& Desarrollador 3D en Tiempo Real especializado en Unity, WebGL y Blender, con experiencia en optimización CAD-2-realidad y creación de visualizaciones técnicas interactivas. Uso Python e IA para automatizar pipelines de producción.”}
\end{quote}




Este enunciado resalta su especialización Unity/WebGL, optimización 3D, visualización técnica e integración de AI/Python. Se puede ajustar según el rol:




\begin{itemize}
\item \textbf{Unity Technical Artist / 3D Developer:} “Technical Artist Unity con enfoque en render en tiempo real y UI interactivas.”  
\item \textbf{Interactive 3D Developer / WebGL Developer:} “Desarrollador 3D interactivo (Unity WebGL) para visualización técnica.”  
\item \textbf{XR/Simulación/Digital Twin:} “Desarrollador XR/Simulación especializado en gemelos digitales para manufactura y drones.”  
\item \textbf{Python/AI Tooling:} “Desarrollador Python \& AI construyendo herramientas de automatización de investigación (e.g. ARA Framework).”  
\end{itemize}




Se recomienda priorizar Unity/WebGL técnicos en industrias de visualización industrial, drones, energía o defensa, antes que juegos puros (donde suele requerirse portafolio artístico más extenso). El candidato puede mencionar videojuegos/tecnología creativa como pasión, pero por estrategia resaltar \textbf{productividad / visualización técnica} (dónde hay más demanda remunerativa). Ejemplo: “Real-Time 3D Developer (Unity) con foco en simulación industrial y visualización técnica, complementado con herramientas de IA para mejorar pipelines.”




La combinación “Technical Artist \& Real-Time 3D Developer” cubre múltiples roles listados (1,2,3,4,5,6,8), mientras que el componente “Python AI Tooling” (punto 11) se presentaría como secundario en CV/LinkedIn por si las oportunidades de visualización no fructifican. La ruta full-stack (12) puede omitirse a favor de especialización técnica, aunque conocimientos web/Python pueden destacarse en sección de habilidades o proyectos. 




\section{Matriz de priorización de roles}




\begin{center}
\scriptsize
\begin{longtable}{>{\raggedright\arraybackslash}p{0.184\linewidth} >{\raggedright\arraybackslash}p{0.184\linewidth} >{\raggedright\arraybackslash}p{0.184\linewidth} >{\raggedright\arraybackslash}p{0.184\linewidth} >{\raggedright\arraybackslash}p{0.184\linewidth}}
\toprule
Rol candidato & Demanda Relativa & Ajuste al Perfil & Salario Potencial (@EE.UU./UE) & Comentarios \\
\midrule
\endfirsthead
\toprule
Rol candidato & Demanda Relativa & Ajuste al Perfil & Salario Potencial (@EE.UU./UE) & Comentarios \\
\midrule
\endhead
\textbf{Unity Technical Artist / 3D Dev} & Alta (videojuegos, tech art) & Muy alta & \textbackslash{}\$80k–\textbackslash{}\$130k+/año & Enfocar portafolio en optimización Blender, shaders, UI interactiva. \\
\textbf{Real-Time 3D / WebGL Developer} & Alta (industrial, simulación, web) & Alta & \textbackslash{}\$70k–\textbackslash{}\$120k/año (XR, simulación) & Prototipos interactivos como TwinSight X500 atraen en industria técnica. \\
\textbf{XR / Simulación / Digital Twin Developer} & Media-alta (AR/VR crece) & Media-alta & \textbackslash{}\$90k–\textbackslash{}\$140k/año & Requiere enfatizar hardware/no obstante; mostrar adaptabilidad de skills. \\
\textbf{Technical Visualization (industrial)} & Media (ingeniería, salud) & Media & \textbackslash{}\$80k–\textbackslash{}\$120k/año & Enfatizar proyectos CAD-3D, visualización de procesos técnicos. \\
\textbf{Game Technical Artist} & Media (videojuegos AAA o indie) & Media-baja & \textbackslash{}\$70k–\textbackslash{}\$130k/año & Alto costo de entrar sin experiencia significativa; mejor pivot a sectores civiles. \\
\textbf{Tools Developer / Pipeline (C\#/Python)} & Baja-media & Media & \textbackslash{}\$60k–\textbackslash{}\$100k/año & Podría sumar valor, pero mercado saturado de dev tools; se usa como skill adicional. \\
\textbf{Python/AI Tooling Developer} & Alta (startups, IA/LLM) & Media & \textbackslash{}\$80k–\textbackslash{}\$140k/año & Mantener proyecto ARA en stack, pero no abandonar el enfoque 3D principal. \\
\textbf{Full-stack Developer (Secundario)} & Media (back-end/front-end) & Media-baja & \textbackslash{}\$70k–\textbackslash{}\$120k/año & Como ruta secundaria si otras fallan; mejorar React/FastAPI podría abrir puertas B2B. \\
\bottomrule
\end{longtable}
\end{center}




\textbf{Conclusión rol-prioridad:} Priorizar \emph{Unity Technical Artist / Real-Time 3D} enfocado en industrias de simulación/visualización (puntos 1,2,3,5,6,8). Incluir \emph{XR/Simulación} (7,8) como roles afines. Posicionamiento secundario en \emph{AI/Python} (11) para roles de herramienta si conviene, pero sin diluir la narrativa principal en 3D. Se descartaría dedicación a full-stack puro (12), más allá de lo básico web, pues aleja del objetivo 3D.




\section{Tabla de priorización de mercados}




\begin{center}
\scriptsize
\begin{longtable}{>{\raggedright\arraybackslash}p{0.184\linewidth} >{\raggedright\arraybackslash}p{0.184\linewidth} >{\raggedright\arraybackslash}p{0.184\linewidth} >{\raggedright\arraybackslash}p{0.184\linewidth} >{\raggedright\arraybackslash}p{0.184\linewidth}}
\toprule
Región/País & Mercado remoto/contratista & Demanda 3D/Unity/Tech Art & Salario USD ≈ (junior–mid) & Comentarios estratégicos \\
\midrule
\endfirsthead
\toprule
Región/País & Mercado remoto/contratista & Demanda 3D/Unity/Tech Art & Salario USD ≈ (junior–mid) & Comentarios estratégicos \\
\midrule
\endhead
\textbf{EE.UU. / Canadá} & Muy alta (tech hubs, estudios juegos, startups) & Alta (videojuegos, AR/VR) & \textbackslash{}\$60k–\textbackslash{}\$130k+/año (jun-senior) & Prioridad alta. Contratos B2B o EOR posibles. Buen pago pero alta competencia. \\
\textbf{Europa Occidental (DE, UK, NL, Nordics)} & Media-Alta (empresas tecnológicas, simulación) & Alta (industria, automotriz, VR) & €40k–€80k/año junior, mid €60k–€100k (≈\textbackslash{}\$45k–\textbackslash{}\$110k) & De alta calidad, aunque posible barrera idiomática (salvo inglés en tech). Evaluar Alemania tras reubicación. UK/Países Bajos muy remotos con inglés. \\
\textbf{Escandinavia (Suecia, Noruega, Finlandia)} & Media (tech y juegos) & Alta (juegos, visualización médica) & Similar Europa (\textasciitilde{}\textbackslash{}\$50k–100k) & Alto costo de vida, buen inglés, sectores VR/medical. Menos remotos. \\
\textbf{Sudeste Asiático} (por evid) & Baja (pocos roles remotos locales, pocos estudios) & Baja & \textbackslash{}\$30k–\textbackslash{}\$50k & Dejar como opción final; salaje bajo relativo. \\
\textbf{LatAm (remoto)}: México, Brasil, Argentina & Alta (nearshore, outsourcing) & Media (outsourcing dev/games) & USD 25k–50k (contratista) & Buen ajuste idiomático (español/inglés). Competencia local media. Salarios menores comparados a EE.UU., pero posibles para Ingresos iniciales. \\
\textbf{España / Portugal} & Media (creciente tech, gaming) & Media (videojuegos, digital twin) & €20k–€40k junior, mid €30k–€60k (\textbackslash{}\$22k–\textbackslash{}\$65k) & Cuidado con burocracia UE. Inglés a veces suficiente; portugués útil a largo plazo. \\
\textbf{Alemania / Europa centro} & Media-Baja (sector manufactura y VR) & Media (ingeniería, simulación) & €30k–€60k (\textbackslash{}\$33k–\textbackslash{}\$66k) & Idioma puede ser barrera; visa familiar al casarse en Alemania abre mercado. \\
\textbf{India / Middle East} & Baja (focos distintos) & Baja & \textbackslash{}\$15k–30k & No estratégico. \\
\textbf{Remoto global (empresas fully-remote)} & Muy Alta & Alta (todas) & variable (\textbackslash{}\$40k–\textbackslash{}\$120k) & Plataformas remotas (GitLab, Automattic) buscan perfiles tecnológicos, aunque no necesariamente 3D. Buena oportunidad freelance / contract. \\
\bottomrule
\end{longtable}
\end{center}




\textbf{Recomendaciones de mercado:} Dar prioridad a \textbf{EE.UU./Canadá} y \textbf{Europa occidental} por remuneración y demanda técnica. También explorar \textbf{LatAm para puestos remotos bilingües} (empresas nearshore y startups) que ofrezcan USD/competitividad. Países de habla inglesa facilitan entrevistas sin dominio alemán. Preparar CV y portafolio adaptado a cada mercado (ej. traducir al inglés, formatos internacionales). Evitar mercados locales de salario bajo (Colombia) y Asia semieconómica. 




\section{Tabla de benchmarking salarial}




Se presenta un cuadro resumen de rangos salariales estimados por mercado y rol, en USD mensuales (equivalente bruto aproximado):




\begin{center}
\scriptsize
\begin{longtable}{>{\raggedright\arraybackslash}p{0.184\linewidth} >{\raggedright\arraybackslash}p{0.184\linewidth} >{\raggedright\arraybackslash}p{0.184\linewidth} >{\raggedright\arraybackslash}p{0.184\linewidth} >{\raggedright\arraybackslash}p{0.184\linewidth}}
\toprule
Rol / Mercado & EE.UU. remoto (USD) & UE remoto (USD) & Contratista B2B Col (USD) & Fuente/Comentarios \\
\midrule
\endfirsthead
\toprule
Rol / Mercado & EE.UU. remoto (USD) & UE remoto (USD) & Contratista B2B Col (USD) & Fuente/Comentarios \\
\midrule
\endhead
\textbf{Unity Technical Artist / Dev} & \textbackslash{}\$6,000–10,000 (mid) & \textbackslash{}\$4,500–8,000 (mid) & \textbackslash{}\$2,000–3,500 & Glassdoor: Tech Artist remoto \textasciitilde{}\$90k\textbackslash{}\textasciitilde{}\textbackslash{}\$7.5k \\
\textbf{XR/AR/VR Developer} & \textbackslash{}\$8,000–12,000 & \textbackslash{}\$5,000–9,000 & \textbackslash{}\$2,500–4,000 & XR dev Glassdoor med \textbackslash{}\$113k \\
\textbf{Technical Visualization} & \textbackslash{}\$6,000–10,000 & \textbackslash{}\$4,000–8,000 & \textbackslash{}\$2,000–3,500 & Similar a above, roles industriales altos. \\
\textbf{3D WebGL/Interactive Dev} & \textbackslash{}\$5,000–9,000 & \textbackslash{}\$4,000–7,000 & \textbackslash{}\$2,000–3,000 & Unity WebGL enfocado a web/móviles, quizá menor que AR/VR. \\
\textbf{Game Technical Artist} & \textbackslash{}\$6,000–12,000 & \textbackslash{}\$5,000–9,000 & \textbackslash{}\$1,800–3,000 & Alta mediana EE.UU., pero baja medio local; campo competitivo. \\
\textbf{Python/AI Tools Dev} & \textbackslash{}\$7,000–12,000 & \textbackslash{}\$5,000–9,000 & \textbackslash{}\$2,500–4,000 & STEM/AI roles IT mediano-alto, en AI startups puede subir. \\
\textbf{Junior Generalist (LatAm)} & \textbackslash{}\$2,500–4,000 & \textbackslash{}\$2,000–3,500 & \textbackslash{}\$1,500–2,500 & Salarios latinos en USD, a veces pagados en dólares. \\
\bottomrule
\end{longtable}
\end{center}




\emph{Fuentes:} Datos de Glassdoor para roles relacionados, reportes ZipRecruiter/Levels.fyi. Los rangos UE son aproximados (Euro a USD). Los valores locales B2B/Col se estiman en base a tarifas freelance en plataformas tech (\textasciitilde{}\textbackslash{}\$15–30/hora según experiencia) y salidas típicas de proyecto (\textasciitilde{}\textbackslash{}\$2000–3500/mes inicial).  




\textbf{Comentarios:} Los valores en EE.UU./UE representan salarios brutos anuales convertidos mensualmente. Los salarios como contratista B2B desde Colombia suelen pactarse en USD (pagos internacionales) y son competitivos para salarios LATAM, pero inferiores al nivel EE.UU. Por ejemplo, un contrato mid-level Unity puede rondar \textbackslash{}\$25–50/hora (USD) en Outsourcing, unos \textbackslash{}\$4k–8k/mes, mientras que un contrato similar en EE.UU. fijo ofrece más.




\section{Análisis de factibilidad de remoto desde Colombia}




Trabajar remoto desde Colombia como contratista (B2B) es viable, pero requiere cumplir normativas:




\begin{itemize}
\item \textbf{Formato legal:} Se utiliza un \emph{Contrato de Prestación de Servicios} a nombre de un trabajador independiente. No hay retención de impuestos por el cliente extranjero, pero el consultor/contratista (Alexander) debe facturar legalmente.
\end{itemize}




\begin{itemize}
\item \textbf{Impuestos:} El contratista presenta \emph{Formulario 210} anual de renta y paga IVA 19\% sobre servicios locales. Sin embargo, los servicios exportados (facturables a empresas extranjeras) están gravados al 0\%. Debe registrarse en RUT (Registro Único Tributario) y emite \textbf{factura electrónica} mediante DIAN (con IVA 0\%). Sin RUT válido, aplican retenciones altas.
\end{itemize}




\begin{itemize}
\item \textbf{Seguridad social:} Como independiente, Alexander debe cotizar autónomamente a salud/pensión (EPS y Colpensiones/otra AFP) y ARL (riesgos laborales) por su cuenta. Los clientes pueden exigir comprobante de ARL vigente. Estos aportes garantizan cobertura.
\end{itemize}




\begin{itemize}
\item \textbf{Pago:} Lo ideal es transferencia bancaria internacional o plataformas como Payoneer/TransferWise (Wise). Se debe declarar correctamente el movimiento de divisas al Banco de la República. Papaya aconseja pagos por wire USD o PayPal (respetando reglas AML).
\end{itemize}




\begin{itemize}
\item \textbf{Otros:} Registrar empresa unipersonal/independiente (“persona natural comerciante”) o al menos RUT. Llevar contabilidad básica (facturas, comprobantes EPS/ARL). Eventualmente considerar EOR (no necesario si se gestiona B2B directo).
\end{itemize}




\textbf{Conclusiones:} Conformar una formalidad básica (RUT, facturación electrónica, ARL/seguridad social) es esencial para acceder a clientes internacionales y evitar penalidades. Se recomienda revisar guías legales: Umbrex y Mismo coinciden en procedimientos de contratación y evita’er “misclasificación”. Se incluirá un checklist: “Registrar RUT, contratar ARL, abrir cuenta COP/USD, habilitar Factura Electrónica DIAN, definir condiciones de pago internacional”.




\section{Estrategia futura UE (Alemania/Portugal)}




\begin{itemize}
\item \textbf{Alemania:} Si la pareja alemana de Alexander logra empleo allí, Alexander puede solicitar visa de reagrupación familiar. Como cónyuge de ciudadano alemán (UE), no se requiere alemán para la visa. Entonces obtendría permiso de residencia/trabajo sin mayores trabas y podrá buscar empleo on-site. Sin embargo, esto depende de los planes reales de la pareja (no garantizado). A medio plazo, si la pareja se traslada, Alexander debe maximizar su inglés, quizá empezar alemán básico para integración. No asumir privilegios alemanes antes de casarse y obtener visa.
\end{itemize}




\begin{itemize}
\item \textbf{Portugal (ciudadanía Sephardí):} Actualmente la ley exige \emph{residencia física 3 años en Portugal} para obtener la ciudadanía por linaje. Si Alexander inicia este proceso, deberá residir en Portugal como mínimo 3 años (puede ser no consecutivos) para poder naturalizar. En el corto plazo (3–6 meses) esto no está disponible. Es un proyecto a largo plazo (≥3 años). No debe mencionarse como “inminente” para reclutadores, solo como meta personal eventual. Conviene mantenerlo en privado en los primeros meses.
\end{itemize}




\begin{itemize}
\item \textbf{Otras opciones UE:} Existen visas de residencia para trabajadores digitales en países como España, Portugal (sin citizenship), pero requieren ingresos mínimos sostenibles (\textasciitilde{}€2000–3000/mes) y a veces prueba de fondos. Podría explorarse en 12–24 meses si tiene experiencia laboral internacional. Por ahora, es secundario.
\end{itemize}




\section{Estrategia de CV}




Se preparará un CV orientado a ATS (un página) en inglés, con claro perfil y bullets cuantificados. Estructura:


\begin{itemize}
\item \textbf{Encabezado:} Nombre, título (p.ej. \emph{“Technical Artist \& Unity Developer”}), localización (Colombia), correo/LinkedIn/GitHub.
\item \textbf{Resumen profesional (2–3 líneas):} Destacar  perfil \emph{“Multimedia Engineer / Technical Artist especializado en Unity WebGL y optimización 3D; apasionado por visualización técnica y producción asistida por IA.”}.
\item \textbf{Habilidades clave:} Lista de Skills en bloque (Unity, C\#, Blender, WebGL, Shader Graph, Python, IA, etc.) con términos ATS.
\item \textbf{Experiencia (proyectos relevantes):} Dado que no hay empleo formal relevante, se listarán proyectos clave como si fueran “freelance/independiente”:
\item \emph{TwinSight X500 (Proyecto de Tesis – Interactive 3D Drone Visualization)}: bullets con métricas (“Optimización de CAD 6.5M→95k triángulos; desarrollo Unity WebGL, UI técnico; evaluaciones SUS 91.88, NASA-TLX 8.69 vs 19.89”). Mostrar liderazgo técnico y alcance.
\item \emph{ARA Framework (Automatización de Investigación Académica, Python)}: bullets (“Desarrollo de sistema de agentes LLM; integraciones LangGraph, scraping de papers; validado parcialmente con prototipos. Presentar como proof-of-concept de AI tooling.”).
\item \emph{Otros proyectos destacados:} (Adaptador: Por completarse en CV bullets: ai-news-aggregator, FitApp-Free, retrato Blender). Cada uno con logros medibles (p.ej. FitApp: “Back-end Flask FastAPI con +100 usuarios activos”; retrato Blender: “Pipeline completo característico (modelado, texturizado, iluminación, render)”).
\end{itemize}




\begin{itemize}
\item \textbf{Educación:}  
\item UNAD – Ingeniero Multimedia (Agosto 2021 – Jul 2026). Indicar mención de \emph{tesis TwinSight X500 / WebGL} (no “ARA Framework”). Resaltar materias de gráficos, software.  
\item UNAL – Ingeniería Electrónica (Marzo 2013 – Nov 2018, 8 semestres completados). Frase honesta: “Studied Electronic Engineering (80\% of curriculum) at Universidad Nacional de Colombia” o “Completed advanced coursework in Electronic Engineering”.  
\item \textbf{Idiomas y herramientas:} Listado de idiomas (Español nativo, Inglés C1). Herramientas (Unity, Blender, Git, etc.).
\item \textbf{Cursos:} Mencionar sin certificar como “Training en pipeline de personajes” o “3D animation” (p.ej. “CG Cookie: High-Fidelity Character Pipeline \& Topology (self-study)”), si es de alguna relevancia.
\end{itemize}




\textbf{Palabras clave ATS:} Unity, Real-Time 3D, WebGL, C\#, Python, Shader Graph, Performance Optimization, XR, Visualization, Technical Artist, Blender, GameDev, Automation, LangGraph, LLM, AI Tools, GitHub, LinkedIn, Portfolio.




Se desarrollarán \textbf{4 variantes de CV} focalizadas:


\begin{enumerate}
\item \emph{Unity Technical Artist / Real-Time 3D}: Enfatizar proyectos TwinSight, Shaders, Blender optimization. Habilidades Unity centradas.
\item \emph{Unity WebGL / Interactive 3D}: Poner TwinSight y habilidades web (URP, WASM, UI Toolkit). Mostrar portafolio web.
\item \emph{XR / Simulación / Digital Twin}: Añadir palabras clave XR/VR, sim. Más énfasis en prototipos industriales, CAD.
\item \emph{Python AI Tools}: Destacar ARA Framework, ai-news, habilidades Python/AI. Dejar TwinSight como secundario, enfocándose en LLM/multiagent (evaluar si combinar o mantener separado del main profile).
\end{enumerate}




\section{Banco de bullets CV}




Ejemplos de viñetas cuantificables para CV (en inglés): 




\begin{itemize}
\item \textbf{TwinSight X500 (Technical Visualization):} 
\item “Developed Unity WebGL prototype for interactive 3D visualization of a Holybro X500 drone; optimized CAD assemblies from 6.5M to 95K triangles for real-time performance.”  
\item “Implemented exploded view, cross-section, selection UI and multiple render modes (realistic, x-ray, blueprint, thermal) using Unity URP and Shader Graph.”  
\item “Orchestrated academic evaluation with 12 participants: achieved SUS=91.9 and NASA-TLX raw=8.69 (3D) vs 19.89 (2D), demonstrating UX improvement.”  
\item “Led asset optimization in Blender (retopology, baking) and pipeline from STEP CAD to low-poly game-ready models.”  
\end{itemize}




\begin{itemize}
\item \textbf{ARA Framework (AI Research Automation):} 
\item “Built a multi-agent Python system (LangGraph, LangChain) for automating literature search and report generation; integrated Semantic Scholar API and browser automation (Playwright/MarkItDown).”  
\item “Designed agents for niche analysis, literature review, architecture drafting and content synthesis; validated with initial mock outputs.”  
\item “Prepared the architecture and prototype pipeline (Redis, Groq LLaMA model) to demonstrate LLM orchestration in research workflows.”  
\end{itemize}




\begin{itemize}
\item \textbf{AI News Aggregator:} 
\item “Developed full-stack AI-driven news platform (FastAPI, React, PostgreSQL, Redis) to crawl and analyze tech news via LLM pipelines.”  
\item “Implemented multi-agent pipeline for news scraping and sentiment analysis; tested MVP with [n] feed integrations.”  
\end{itemize}




\begin{itemize}
\item \textbf{FitApp-Free (Fitness App):} 
\item “Prototyped mobile fitness app backend (Flask/FastAPI) for calisthenics training; features user progress tracking, injury alerts, and adaptive routines.”  
\item “Configured PostgreSQL data models for workouts and progress; integrated React frontend scaffolding.”  
\end{itemize}




\begin{itemize}
\item \textbf{Blender Retrato Hiperrealista:} 
\item “Modelled and rendered a high-fidelity human portrait in Blender following CG Cookie pipeline, demonstrating advanced topology, sculpting, texturing and lighting skills.”  
\item “Produced wireframe topology charts, UV layout and shading breakdowns to show technical process; published case study on ArtStation.”  
\end{itemize}




\begin{itemize}
\item \textbf{Generales:} 
\item “Published and maintain GitHub repos for portfolio; cleaned/organized code (private→public) and curated README for recruiters.”  
\item “Launched personal portfolio site (HTML/CSS/JS) with project demos, CV, and case studies to showcase work to international recruiters.”  
\end{itemize}




Estos bullets deben adaptarse a cada versión del CV, usando verbos de acción y datos numéricos (por ejemplo, porcentaje de reducción de polígonos o métricas de usabilidad) cuando sea posible. Evitar el término “AI” sin contexto: en entrevistas se enfatizará que fue \emph{asistencia} y su rol fue diseño/programación.




\section{Estrategia de LinkedIn}




\textbf{Foto de perfil:} Profesional, fondo neutral, rostro claramente visible. Si es posible, una imagen "técnica" (p.ej. frente a modelado 3D) pero preferiblemente primer plano amistoso/profesional. Evitar selfies informales.




\textbf{Titular (“Headline”)}: Enfocado en palabras clave técnicas. Ejemplos:


\begin{itemize}
\item \emph{Unity Technical Artist \& Tools Dev} – Unity, WebGL, C\#.
\item \emph{Real-Time 3D Developer (Unity/WebGL) – Technical Visualization}.
\item \emph{Technical Artist | Unity/WebGL | 3D Optimization}.
\item \emph{Multimedia Engineer / Unity Developer – IA \& Python Tooling}.
\item \emph{Technical Artist – Blender, CAD-to-RealTime, AI-Driven Pipelines}.
\end{itemize}




Se darán \textasciitilde{}15 opciones similares y se elegirá la más concisa/potente. Ejemplo final: \textbf{“Technical Artist \& Real-Time 3D Developer (Unity, C\#, WebGL) – 3D Visualization \& AI Tooling”}.




\textbf{Extracto/“About” (EN):} 2–3 párrafos cortos. Debe ser equilibrado entre profesional y personal, “premium”. Mencionar punto fuerte y metas: ej. “Multimedia Engineer with Unity \& Blender expertise, passionate about transforming CAD data into interactive 3D experiences. Completed thesis “TwinSight X500”: Unity WebGL drone visualization with 6.5M→95K triangle optimization. Seeking remote roles in Technical Art/Visualization to leverage my background in realtime graphics and AI-assisted tools. Spanish native, English fluent, open to global opportunities.” En español: síntesis breve de misión (“ayudar a ingenieros a entender sus diseños 3D, construyendo interfaces web immersivas…”), pero preferiblemente en inglés para coherencia internacional. No repetir CV/portafolio, sino contar una historia concisa del perfil, indicando disponibilidad remota (“based in Colombia, open to remote or EU relocation”) y aspiraciones (ej. “dynamics roles requiring 3D dev \& creative tech”). Evitar tono modesto, ir seguro de logros.




\textbf{Experiencia / Proyectos en LinkedIn:} Reescribir “Independent Technical Artist”:


\begin{itemize}
\item Título: “Freelance Technical Artist / Unity Developer” (Mar 2024–Presente).
\item Descripción: “Desarrollo herramientas automatizadas y optimización de activos 3D; integración de assets URP; creación de shaders (dissolve, hologramas, etc.). Gestión de pipeline Blender→Unity para proyectos de visualización. Destacar “TwinSight X500 (Tesis): Unity WebGL para visualización dron, optimización 99\% de malla y UI técnica”, y “ARA Framework: pipeline IA para investigación académica usando LangGraph”. Incluir métricas breves.
\end{itemize}




\textbf{Educación:} 


\begin{itemize}
\item UNAD Multimedia (2021–2026). Descripción: “Engineering multidisciplinary (3D, software, matemática) con énfasis en gráficos y realidad virtual; tesis final: \emph{TwinSight X500, prototipo web 3D interactivo}.”  
\item UNAL Electrónica (2013–2018, 80\% curr.). Descripción: “Advanced coursework in Electronic Engineering (electronics, systems modeling, optimization). Systems simulation, física, lógica digital, etc.”  
\end{itemize}


Voluntariado ARCADE – “Student Researcher (Feb 2025–Presente): investigación en visualización 3D y XR en entornos educativos.”




\textbf{Sección “Proyectos”:} Incluir “TwinSight X500 – WebGL Drone Assembly” (descripción breve, e.g. “Unity WebGL app for interactive drone assembly visualization (optimización CAD, UI, cross-sections) – \emph{GitHub/Proyecto}”). “ARA Framework – LLM Research Tool”, “Blender Portrait – Character pipeline” (ligar a ArtStation/GitHub).




\textbf{Destacados (“Featured”):} Ordenar: Portafolio/Media principal (video demo o imágenes TwinSight), GitHub, ArtStation, publicaciones relevantes. TwinSight debe ser lo primero, luego ARA (si público), luego otros proyectos.




\textbf{Skills (Aptitudes):} Priorizar 7–10 más relevantes: \emph{Unity, 3D Graphics, C\#, WebGL, Technical Art, Python, Git, 3D Optimization, VR/AR, Shader Graph, Blender}. Reordenar para que aparezcan primeros. Añadir “B2B Contracts” o “Remote Work” en idiomas o intereses si aplica. El apartado de “Languages” ya está bien (Spanish, English).




\textbf{Perfil/Apariencia:} Banner de LinkedIn sugerido: Una composición profesional de renderizado 3D de un dron o escenas técnicas (ej. captura del TwinSight X500) con texto superpuesto como “Unity | Technical Artist | 3D Developer”. Foto profesional (no cara al sol, etc). 



\textbf{Actividad:} Publicar contenido relevante periódicamente: estudios propios (por ejemplo, avances de TwinSight X500, infografías), artículos cortos (ej. “Optimizing CAD models for real-time: mi experiencia con drone X500”), re-publicar noticias de XR/técnica. Comentar en posts de Unity/Blender, participar en grupos de “Tech Art” o “Unity Developer” en LinkedIn, ArtStation. Esto mejora visibilidad.




\textbf{Estrategia de palabras clave (recruiters):} Incluir en resumen/títulos variantes bilingües: \emph{Technical Artist}, \emph{Desarrollador 3D}, \emph{Unity Developer}, \emph{WebGL Engineer}, \emph{XR Developer}, \emph{Visualization Engineer}, \emph{Herramientas de Producción IA}, etc. El algoritmo de LinkedIn: asegurarse de mencionar “remote” y “Colombia” (o LatAm) en la sección “Información” para aparecer en búsquedas como “Unity developer remote Latin America”. Ejemplo booleano para reclutadores: 


\begin{mdcode}
Bloque de codigo
("(Unity Technical Artist" OR "3D Developer" OR "Interactive 3D Developer" OR "WebGL Developer" OR "XR Developer" OR "Digital Twin Developer") AND (remote OR Colombia OR LATAM OR remoto) 
\end{mdcode}


Esto ayuda a que reclutadores que buscan perfiles 3D/Unity/remote nos encuentren.




\textbf{Disponibilidad geográfica en perfil:} Mencionar “Based in Colombia (Remote only) / Willing to relocate to EU (esp. Germany/Portugal) if visa provided”, en la sección de “Intereses” o “Resumen” sin sonar pretencioso: por ejemplo, “Open to remote roles globally; Long-term EU mobility (via Portuguese citizenship in process)”. Esto semáforo: nada de afirmar “Ciudadano UE inminente” sin cumplir requisitos.




\textbf{Otras:} Evitar inconsistencia: corregir en LinkedIn actual que pone “Tesis: ARA Framework” si en realidad fue TwinSight. Usar sincronización: el título “Technical Artist \& Tools Developer” está bien, acentuar “Optimization, WebGL” en la sección. Asegurarse de quitar proyectos irrelevantes o sacar el portafolio duplicado en Educación. 




\section{Estrategia de GitHub}




\begin{itemize}
\item \textbf{Repositorios clave:} Definir 3–5 repos “estelares”. Poner en \textbf{pinned} (Bio de GitHub):
\end{itemize}

\begin{enumerate}
\item \emph{WebGL-Thesis-Proposal} (TwinSight X500) – principal. Limpieza: añadir README explicativo conciso, capturas, licencia. Incluir \emph{Demo deploy} link (si posible host en GitHub Pages/WebGL).  
\item \emph{ARA-Framework} – si se decide público. README claro: “Python LLM Orchestration for research”. Señalar que es proto/incompleto.  
\item \emph{ai-news-aggregator} – si finaliza mínimamente, poner con demo/README si lo publica.  
\item \emph{3D Character (Blender Portrait)} – usar ArtStation como medio y GitHub solo para wireframes etc.  
\item \emph{GitHub-intro / Portfolio repos} – eliminar, archivar o combinar.
\end{enumerate}




\begin{itemize}
\item \textbf{Público vs privado:} Convertir en público los proyectos más fuertes (TwinSight, proyectos de AI/pipeline). Mantener privado o archivar códigos sin valor reclutable (ej. proyectos obsoletos, antiguos), o señalarlos como \textbf{portfolio antiguo}. Poner notas README aclarando estado (beta, no mantenimiento).  
\item \textbf{README}: En cada repo, inicializar con un “recruiter-friendly” README: breve descripción, tecnología, rol propio (ej. “Lead developer and integrator”), demos/video embed, métricas clave (TiPs). Evitar jergas de nicho y destacar resultados/uso.  
\item \textbf{GitHub Pages}: Crear página portfolio vinculado a delarge95.github.io (por ejemplo con los proyectos destacados).  
\item \textbf{Etiquetado:} Crear etiquetas y temas relevantes (Unity, C\#, WebGL, Python).  
\item \textbf{Actividad:} Contribuir a repos públicos/iniciativas Unity o CG (por visibilidad). Documentar algunos issues/respuestas en comunidades (StackOverflow).  
\item \textbf{Código limpio:} Borrar demos privados con información sensible, comprimir assets, etiquetar commits.  
\item \textbf{IA Asistida:} Si se ha usado ChatGPT en código, no ocultarlo pero dejar claro que el trabajo final fue suyo. En README, se puede mencionar “Auxiliary use of AI tools for code scaffolding; all design and logic by author.”
\end{itemize}




\section{Estrategia de portafolio}




\begin{itemize}
\item \textbf{Sitio portfolio (MVP):} Construir un sitio web sencillo (hosteable en GitHub Pages) que enlace a proyectos: \emph{TwinSight X500} como caso estrella, con video demo breve (60–90s), capturas, descripción. Incluir \emph{ARA Framework}, \emph{AI-news} solo si se considera relevante, pero indicarlo como “herramientas internas” (secundario). \emph{Blender Portrait} en sección Art/Graphics. Cada proyecto con resumen, rol, tecnologías, logros cuantificables (tomados de CV bullets).
\end{itemize}




\begin{itemize}
\item \textbf{Case studies:} Para cada proyecto principal (TwinSight, ARA, Blender), preparar página con:
\item \textbf{Problema/desafío}: contexto (p.ej. necesidad de entender dron industrial en 3D).  
\item \textbf{Solución}: enfoque técnico (UX, tecnologías usadas).  
\item \textbf{Rol personal}: qué contribuciones propias se hicieron (enfatizar partes críticas).  
\item \textbf{Resultado}: métricas o estado (objetivos logrados, validaciones).  
\item \textbf{Lecciones / Futuro}: mejoras pendientes o aprendizajes (muestra reflexión profesional).  
\end{itemize}

  Incluir screenshots (wireframe, editor Unity, UI, explosion, heatmap de performance, etc.), gráficos (reducción polígonos, resultados de SUS/TLX).




\begin{itemize}
\item \textbf{ArtStation:} Complementario. Cargar retrato Blender. Estructura técnica: imágenes del modelo wireframe/texture/bake, explicaciones breves en cada. Conectar a \textbf{Technical Art}: mostrar fluidez con Blender (no solo arte final). Possibly mostrar antes/después wireframe→render.
\end{itemize}




\begin{itemize}
\item \textbf{PDF/Deck:} Tener un PDF unificado de portfolio para enviar por email. Que incluya resumen de habilidades, proyectos clave y enlaces. A menudo reclutadores aprecian un “one-pager” visual. Mantener en inglés.
\end{itemize}




\begin{itemize}
\item \textbf{Demo reel/video:} Un video de 60–120s para TwinSight X500: 
\item \emph{Guión}: abrir con contexto (“App to explore drone assembly in 3D”), mostrar interfaz UI seleccionable (click en partes), demonstration of explosion mode, cross-section, x-ray etc. Resaltar performance en móvil/web. Concluir con “Ver proyecto completo en GitHub/portafolio.”
\item \emph{Shots}: listamos escenas cortas: vista general del dron (rotando), click part (muestra panel info), explosion (animación separación de partes), modos visuales (cambiar a x-ray, thermal, blueprint), dispositivo móvil (mostrar que corre en web mobile).  
\item Añadir voz o texto superpuesto breve: “Interactive 3D visualization; CAD→real-time optimization; implemented in Unity WebGL.”
\end{itemize}




\section{Company target list (ejemplos destacados)}




Generar lista de \textasciitilde{}100 posibles empleadores/agencias, priorizados por fit. Ejemplos por categoría (\textbf{resume}):




\begin{itemize}
\item \textbf{Videojuegos AAA/indie:} Ubisoft (FR/Global) – Technical Artist (remoto USA/FR), Microsoft Xbox Game Studios (US), Insomniac Games (US), Bungie (US), Cloud Imperium Games (Star Citizen), Unity Technologies (US), Epic Games (US), Tencent/Nexon (Asia).  
\item \textbf{Realidad Virtual / Simulación:} Meta (US, XR Research), Microsoft (Hololens), NVIDIA (Omniverse), Varjo (FI), Niantic (US), Magic Leap (US), Unity (iluminación).  
\item \textbf{Defensa/Espacio:} Boeing (US), Lockheed Martin (US), Northrop Grumman (US), Raytheon (US), Airbus Defence (EU), DLR (Alemania).  
\item \textbf{Industria / CAD / Automoción:} Siemens Digital Industries (DE), Dassault Systèmes (FR), ANSYS (US), Bentley Systems (US/UK), AutoDesk (US). Roles en visualización o digital twin.  
\item \textbf{Aeroespacial/Drones:} NASA (JPL/GSFC), Cessna/Textron, SpaceX (Motion Graphics, simulación), Holybro (EE.UU.).  
\item \textbf{Visualización médica/educativa:} Siemens Healthineers (DE), GE Healthcare (US), Philips (NL).  
\item \textbf{Agrícola / Energía:} John Deere (US), Schlumberger (US, AR/VR oilfield), ABB (Suiza/SE), Schneider Electric (FR).  
\item \textbf{Plataformas / Consultoría tecnológica:} Globant (ARG, oficinas en COL), Endava, Wipro, Accenture (gamedev/tech art divisions).  
\item \textbf{Startups creativas / AI:} MagicLeap, FormLabs, Cranial (AI demo), VR education startups.  
\item \textbf{Outsourcing LATAM:} UruIT (Col), PSL Corp (Col), Arkano (UY), Demium (ARG), Levvel (adapt tech in Colombia).  
\item \textbf{Gobierno/Educación:} Labs de VR en universidades (p.ej. ETH Zurich simulation lab, Fraunhofer IGD), proyectos financiamiento europeo.  
\end{itemize}




Para cada empresa, idealmente se anotarían \emph{remote likelihood} (p.ej. Unity/Epic/KLaber son fully remote friendly), salarios estimados, roles a aplicar (“Technical Artist”, “3D Dev”), idioma principal, etc. Esto requiere investigación puntual por empresa (carreras, Glassdoor reviews). En esta respuesta no es posible listar 100 con detalle citado, pero se recomienda al candidato crear una hoja de cálculo con criterios: cercanía industrial, historial de contratación remota, culturas de startup, etc.




\section{Bolsas de empleo y comunidades}




\begin{itemize}
\item \textbf{General / remotas:} LinkedIn (filtros remotos Unity), Indeed, Glassdoor, ZipRecruiter. Plataformas de empleo remoto: We Work Remotely, RemoteOK, Working Nomads, FlexJobs.  
\item \textbf{Tech art / Unity:} Unity Connect, Polycount Jobs (foros), ArtStation Jobs, Gamedev.net, CGSociety.  
\item \textbf{XR / AR/VR:} XR Bootcamp job board, Augmented Reality developer Slack, arvr.slack.com, XR Dive.  
\item \textbf{Digital twin / Industrial:} LinkedIn Grupos (Digital Twin Forum), Big Data \& IoT Slack, Industry forums (Product Visualization Discords).  
\item \textbf{Freelance:} Upwork, Freelancer, Toptal (temporales solo), Gigster para devs, Cord, Malt (Europa).  
\item \textbf{LATAM específicas:} elempleo.com, compuTrabajo, Workana; comunidades BogotáTech, Medellín GDC; Tech recruitment firms nearshore (MuleSoft LATAM, etc).  
\item \textbf{Comunidades/Discord/Slack:} GameDev Underground (Discord), Unreal Engine Slack, Unity Developer Groups, Blender Artists Forum. LinkedIn grupos: “Technical Artists World”, “Unity Developers”, “XR/AR Professionals”.
\end{itemize}




\section{Sistema de aplicación y métricas}




\begin{itemize}
\item \textbf{Rutina semanal:} Dedicar 2–3 horas diarias (20–25 semanales) a aplicaciones/alcance. Objetivo \textasciitilde{}5-10 aplicaciones semanales (calidad > cantidad).  
\item \textbf{A/B/C categorías:} Categorizar vacantes:  
\item \emph{A (alta probabilidad)}: Experiencia ideal, fit total.  
\item \emph{B (medio)}: Falta experiencia formal, pero habilidades clave están.  
\item \emph{C (difícil)}: Roles severos (senior AAA, necesitan portafolio top).  
\end{itemize}

  Priorizar A/B, aspirar a 50\% customización personal.  

\begin{itemize}
\item \textbf{Customización:} Adaptar CV+cover letter breve (3-4 líneas) a descripción del puesto: mencionar tecnologías clave del anuncio, ej. “Su oferta pide Unity + WebGL + C\#: he liderado X proyecto con esas stack.”  
\item \textbf{Herramientas:} Llevar un tracker en Notion/Google Sheets con columnas: Empresa, Rol, Fuente (link), Contacto (link perfil reclutador), Fecha aplica, Estado, Notas (adaptaciones, fecha follow-up). Analizar métricas mensuales (\# apps, \# entrevistas, \% respuestas).  
\item \textbf{Outreach y seguimiento:} Conectar en LinkedIn con reclutadores y leads (ver scripts abajo). Seguir con 1–2 mensajes de seguimiento cada 10 días después de aplicar si no contestan.  
\item \textbf{Rechazos:} Establecer criterio: si oferta < USD 1.5k/mo neto o condiciones contractuales inestables (sin contrato claro), declinar con cortesía.  
\item \textbf{Balance:} Semana activa de aplicaciones + networking + mejora portafolio. No saturarse, mantener aprendizaje (cursos/idioma) paralelamente.
\end{itemize}




\section{Plan 30/60/90 días (tentativo)}




\textbf{Días 1–30:}  


\begin{itemize}
\item Finalizar defensa de tesis.  
\item \textbf{CV y LinkedIn:} Crear base CV en inglés y variante específico Unity; reescribir LinkedIn según arriba; optimizar GitHub (limpieza repos, public repos finalizados).  
\item \textbf{Portafolio:} Lanzar sitio web simple (GitHub Pages) con TwinSight X500 y Avatar. Crear demo video breve TwinSight.  
\item \textbf{Redes:} Aumentar contactos LinkedIn de industrias target; enviar solicitudes con mensaje personalizado (ver scripts). Unirme a foros/Discord relevantes.  
\item \textbf{Setup administrativo:} Obtener RUT, Afiliaciones EPS/ARL, abrir cuenta COP+USD y vincular Payoneer/Wise. Configurar facturación (Factura electrónica DIAN).  
\item \textbf{Busqueda:} Preparar lista de empresas (+20 prioritarias), reclutadores de Unity/art tech, grupos de Slack. Postular a ofertas iniciales fáciles (revisión de portales/remotos).  
\item \textbf{Tracker:} Establecer en Notion/Sheets y revisar avances semanalmente.  
\item \textbf{Aprendizaje:} Comenzar curso o proyecto corto para cubrir gap (p.ej. React básico o optimización Unity).
\end{itemize}




\textbf{Días 31–60:}  


\begin{itemize}
\item \textbf{Aplicaciones:} Aumentar número (10–15 por semana), aplicar a roles fuera de portafolio inicial. Personalizar cada solicitud.  
\item \textbf{Portafolio/Proyectos:} Subir mejoras en proyectos (TwinSight, ARA). Decidir si publicar ARA (limpieza final). Producir un caso de estudio corto para blog/LinkedIn sobre TwinSight o AI en CV.  
\item \textbf{Networking:} Realizar contactos directos con empresas (LinkedIn InMails, email), asistir a meetups virtuales/bootcamps de tech art. Pedir recomendaciones de profesores/mentores (script profesor).  
\item \textbf{Preparación entrevista:} Practicar preguntas comunes (Unity, gráficos, optimización). Preparar “pitch” de proyectos (TwinSight, ARA). Ensayar respuesta “¿Cuánto código fue tuyo vs IA?” (p.ej. “lideré la arquitectura e integración; la IA ayudó en esqueleto de código, pero implementé y probé cada módulo”).
\item \textbf{Ajustes:} Recoger feedback de primeras entrevistas (si las hay) y mejorar CV/gabs. Ajustar estrategia de rol si es necesario (ej. mayor foco Python si Unity se resiste).  
\item \textbf{Documentación:} Crear plantillas de cover letters, CV, extracts LinkedIn por sector. Empezar a enviar saldos para Alemania/Portugal (idioma o doc, si aplica a largo plazo).
\end{itemize}




\textbf{Días 61–90:}  


\begin{itemize}
\item \textbf{Escalar outreach:} Aplicar a vacantes más ambiciosas (niveles B/C), incluyendo empresas top (puede usar referidos). Hacer seguimiento intensivo a contactos creados.  
\item \textbf{Entrevistas:} Idealmente tener entrevistas técnicas y de cultura. Prepararse con portfolio digital (video, slides).  
\item \textbf{Negociación:} Investigar salarios específicos de ofertas obtenidas. Practicar respuestas a “salario esperado” (basado en research de salario medio, rango deseado, ajustes por remoto).  
\item \textbf{Ofertas:} Evaluar propuestas, priorizar aquellas con mejor remuneración y potencial. Negociar en base a rangos del mercado. Rechazar ofertas muy lejos de objetivo sin comprometer esperanza.  
\item \textbf{Capacitación continua:} Dedicarse a gap skills (React, Three.js, Shader Graph avanzado) según feedback de entrevistas.  
\item \textbf{Preparación EU:} Si en aparece opción alemana/portuguesa, planificar documentos. Mantener opciones abiertas.
\end{itemize}




\section{Scripts de acercamiento (Inglés y Español)}




\emph{(Plantillas básicas; personalizar con nombres/motivos específicos.)}




\begin{itemize}
\item \textbf{Reclutador / Conexión LinkedIn (EN):}  
\end{itemize}

  “Hi [Nombre], I’m Alexander Woodcock, a Unity Technical Artist based in Colombia. I found your profile when searching for [role] opportunities at [Compañía]. With expertise in Unity/WebGL and CAD-to-3D visualization (see my TwinSight X500 thesis), I’m very interested in [empresa/puesto]. Could we connect? I’d love to discuss possible fits. Thank you!”  




\begin{itemize}
\item \textbf{Líder técnico de Tech Art (EN):}  
\end{itemize}

  “Hello [Nombre], I’m a Technical Artist (Unity/WebGL) working on an interactive drone assembly prototype (see GitHub WebGL-Thesis-Proposal). I admire [empresa]’s work in [campo XR/simulación/etc.]. I’d appreciate 15 minutes to share how my Unity/C\# \& optimization skills could support your team’s projects.”  




\begin{itemize}
\item \textbf{Fundador estudio creativo (ES):}  
\end{itemize}

  “Hola [Nombre], soy Alexander Woodcock, Ingeniero Multimedia especializado en Unity y visualización 3D (ej. prototipo web de dron en mi tesis). He visto el trabajo de [empresa] en [campo relevante] y me interesa sumarme como Technical Artist/Desarrollador 3D. ¿Podríamos conversar sobre posibles colaboraciones remotas? Adjunto mi portafolio.”  




\begin{itemize}
\item \textbf{Director técnico Unity (EN):}  
\end{itemize}

  “Hi [Nombre], as a Senior Unity Dev at [empresa], you might be seeking strong Technical Artists. I’m a passionate Unity/WebGL developer (8+ semesters engineering) with a portfolio in 3D optimization (6.5M→95K polys). Could we discuss any opportunities at [empresa]? I’m targeting remote or contractor roles.”  




\begin{itemize}
\item \textbf{Empresa XR/Simulación (EN):}  
\end{itemize}

  “Good day, [Nombre]. I’m contacting you about potential roles in XR/simulation at [empresa]. My background includes Unity WebGL, C\# tools dev, and AI-assisted pipelines. At [proyecto], I built a 3D drone visualization with real-time interactivity. Would love to contribute to [empresa]’s simulation projects.”  




\begin{itemize}
\item \textbf{Reclutador videojuego (ES):}  
\end{itemize}

  “Hola [Nombre], actualmente freelance en Unity/C\# con experiencia en herramientas de optimización 3D. Vi que [empresa] busca Technical Artist para juegos en Unity. Aunque mi pasión es también la simulación, tengo portafolio con shaders y UI Unity. Me encantaría explorar oportunidades de contrato remoto.”  




\begin{itemize}
\item \textbf{Follow-up Postulación (EN):}  
\end{itemize}

  “Hello [Nombre], I wanted to follow up on my application for [puesto] submitted on [fecha]. I remain very interested in this role. My Unity and 3D expertise (TwinSight prototype) align well with the job requirements. Please let me know if you need any further info. Thanks again for considering me.”  




\begin{itemize}
\item \textbf{Aplicación especulativa (EN):}  
\end{itemize}

  “Hi [Nombre], I’m Alexander, a Unity Technical Artist in Colombia. I’m reaching out speculatively as I admire [empresa]’s work. I have experience in real-time 3D visualization and AI tool development. If there are upcoming Unity/3D positions (even contractor), I’d be keen to apply. Thank you for your time.”  




\begin{itemize}
\item \textbf{Solicitud de referencia (EN):}  
\end{itemize}

  “Hi [Nombre], I’m looking to expand my network in the [XR/Technical Art] field. As someone respected in the community, could you kindly refer me to any roles or contacts? I’d be glad to support or informally discuss how my skills (Unity, WebGL, Python) could fit.”  




\begin{itemize}
\item \textbf{Profesor / Asesor (ES):}  
\end{itemize}

  “Estimado [Profesor], le escribo para solicitarle una recomendación/introducción a contactos en la industria 3D. Como exalumno de [UNAD/UNAL], apreciaría su orientación sobre cómo aprovechar mis proyectos (p.ej. tesis TwinSight X500) para oportunidades internacionales.”  




\begin{itemize}
\item \textbf{Cliente freelance anterior (ES):}  
\end{itemize}

  “Hola [Nombre], espero esté bien. Quería pedirle un testimonio breve sobre el proyecto que hicimos juntos (optimización web). Estoy enfocando mi portafolio en trabajos 3D y un comentario suyo sería valioso.”  




\begin{itemize}
\item \textbf{Expectativas salariales (EN):}  
\end{itemize}

  “Based on my research and the responsibilities of this position, I expect a salary in the range of [ej. \textbackslash{}\$70k–\textbackslash{}\$80k] annually. Of course, I’m open to discussion depending on the overall package.” (Luego alinear con objetivos: \textbackslash{}\$1,500->\$18k mínimo, según rol; \textbackslash{}\$3k->\$36k; \textbackslash{}\$6k->\$72k)  




\begin{itemize}
\item \textbf{Disponibilidad remota desde Colombia (ES):}  
\end{itemize}

  “Disponibilidad completa remota desde Colombia; constituiría mi condición habitual de trabajo. Cuento con RUT y facturación oficial. Estoy abierto a contratos directos o B2B.”  




\begin{itemize}
\item \textbf{Configuración contractor (ES):}  
\end{itemize}

  “Como contratista independiente, puedo facturar legalmente servicios exportados (IVA 0\%) y cuento con RUT. Prefiero pagos vía transferencia internacional (Wise/PayPal). ¿La empresa utiliza EOR/Deel o prefiere un contrato directo? Estoy preparado para acordar términos transparentes.”  




\begin{itemize}
\item \textbf{Movilidad UE futura (EN):}  
\end{itemize}

  “I plan to pursue Portuguese citizenship (ancestry) in the next 2–3 years, which may facilitate EU relocation. For now, I seek fully remote or international contractor roles that don’t require current EU work authorization.”  




\begin{itemize}
\item \textbf{Preguntas beneficios/contrato (EN):}  
\end{itemize}

  “Could you clarify if the position is contractor or W2? And what benefits or payment frequency are offered? Also, is there an expected full-time schedule (hours) for this contractor role?”  




\begin{itemize}
\item \textbf{Rechazo oferta inestable (EN):}  
\end{itemize}

  “Thank you for the offer. Unfortunately, the compensation (USD X/month) is below my minimum requirement (USD 1500). Additionally, the role’s scope and contract terms are not yet clear. I prefer stable, well-defined positions. I appreciate your time and wish you the best.”  




\begin{itemize}
\item \textbf{Ausencia de experiencia formal (EN):}  
\end{itemize}

  “While I lack formal studio experience, my portfolio projects (ex: TwinSight X500) demonstrate practical Unity and 3D pipeline skills. I am a quick learner with strong fundamentals (ING. Electrónica 80\% cursado) ready to apply in a professional context.”  




\begin{itemize}
\item \textbf{Uso de IA en proyectos (EN):}  
\end{itemize}

  “I used AI tools as coding assistants (to draft boilerplate) but personally architected, integrated, and debugged all code. My contributions include the core algorithms and systems design. For example, I guided the AI generation and refined every script.”  




\begin{itemize}
\item \textbf{“¿Por qué contratarlo?” (ES):}  
\end{itemize}

  “Poseo un perfil único que combina habilidades técnicas (Unity, C\#, optimización 3D, Python) con experiencia en proyectos end-to-end. Mi tesis TwinSight validó con métricas la mejora en usabilidad. Aunque sin trayectoria empresarial clásica, ya demostré liderazgo técnico y creatividad aplicada. Estoy motivado y preparado para agregar valor real desde el día uno.”  




\section{Preparación para entrevistas y negociación}




\textbf{Entrevistas técnicas:} Esperar preguntas sobre Unity/WebGL (cómo optimizar performance, cómo funcionan shaders, ejemplos de pipelines). Reforzar bases de gráficos, estructuras de datos 3D, experiencia con UI Toolkit. Practicar resolución de problemas en vivo (p.ej. pseudocódigo). En cuanto a IA/Scripted: decir que integra AI, pero no se espera que copie conocimiento exacto. 




\textbf{Preguntas comunes:} “Háblame de un proyecto en que lideraste [TwinSight]. ¿Cómo resolviste X reto de optimización? ¿Cómo calibraste el sistema UX?”; “¿Cómo explicarías tu uso de IA en el desarrollo?”; “Expectativas salariales/visado.”; “¿Por qué Unity vs Unreal?”; “¿Cómo te mantienes actualizado en técnicas de render?”. Preparar respuestas basadas en hechos (mencionar métricas de TwinSight, frameworks).




\textbf{Negociación salarial:} Llevar rangos de referencia (ej. “Unity Tech Artist senior en EEUU \textasciitilde{} \textbackslash{}\$100k”), comunicar ambición pero flexibilidad. Empezar por arriba del mínimo aspirado (\textasciitilde{}\textbackslash{}\$2k en Colombia convertidos) y ajustar. Tomar en cuenta beneficios adicionales (equity, remote stipend). Discutir modalidad (contratista vs empleado) y cobertura de impuestos (explicar cómo tributará en Colombia, esperar no deducir impuestoms del pago).




\textbf{Checklist de oferta:} Además de salario, verificar:


\begin{itemize}
\item Tipo de contrato (W2, 1099, EOR).
\item Días vacacionales, flexibilidad horario.
\item Cobertura de herramientas/licencias.
\item Responsabilidades vs título (nivel real del puesto).
\item Opciones de crecimiento (ojos a 12–24 meses).
\end{itemize}




\section{Estrategias de búsqueda final}




\begin{itemize}
\item \textbf{Conservadora:} Enfocarse en ofertas seguras remotas en empresas medianas (ej. outsourcing game studios, consultoras tecnológicas) con salarios modestos pero estables. Prepararse con CV pulido para entrevistas típicas.  
\item \textbf{Equilibrada:} Mezclar solicitudes a startups internacionales remotas de nicho (visualización, IA) y empresas establecidas (Unity, Autodesk). Usar portafolio activo, EOR/Deel si es necesario para asegurar pago. Mantener networking activo y preparación de habilidades (tal vez Inglés perfeccionado).  
\item \textbf{Agresiva:} Target empresas grandes y elevadas (Google AR, Meta, Boeing VR, etc), networkear directamente con altos cargos, presionarse a aplicar a \textasciitilde{}20 empresas top. Mejorar más habilidades (Shader Graph avanzado, Three.js). Uso intensivo de LinkedIn y referidos.  
\end{itemize}




\textbf{Recomendación:} Combinación equilibrada: aplicar tanto en puestos “probables” (nearshore, mediana tech, \textbackslash{}\$36k/año nivel) como “aspiracionales” (grandes, \textbackslash{}\$72k/año). Mientras se gana confianza y experiencia de entrevistas, ir escalando metas. Balancear búsquedas según prioridades salariales y viabilidad legal.  




\section{Brechas de habilidades (ROI)}




Priorizadas por impacto en empleabilidad:




\begin{itemize}
\item \textbf{Unity Performance \& Profiling:} Conocer Profiler, optimizar GPU/CPU (milestones visuales).
\item \textbf{Unity WebGL específica:} Build/WebGL size limits, compression, Addressables para RAM.
\item \textbf{Shaders (HLSL avanzado):} Más allá de Shader Graph, escribir shaders manualmente.
\item \textbf{Three.js/WebGPU:} Diversificar a Web 3D (por si aplican a roles WebGL puro).
\item \textbf{React y full-stack básico:} Mejorar React en ARA/ai-news, entender REST/DB.
\item \textbf{XR básico:} Fundamentos AR/VR (tracking, SteamVR/ARKit).
\item \textbf{Estructura C\#:} Arquitectura limpia (MVVM, patterns) para soporte a escalas grandes.
\item \textbf{FastAPI/Backend:} Pulir FastAPI para la ruta Python.
\item \textbf{LangGraph/LLM Aplicaciones:} Entender mejor orquestación multi-agent (LangChain nuevo).
\item \textbf{Blender polishing:} Esculpir/retopología/nody vs linear workflow.
\item \textbf{Videography:} Edición básica (Premiere), storyboarding para demo.
\item \textbf{Escritura técnica:} Redactar documentación de proyectos y CV fuerte.
\item \textbf{Idiomas:} inglés (presentations), alemán/portugués básico (entrevistas introductorias).
\item \textbf{Gestión de proyectos Agile:} Familiarizarse con JIRA, Scrum (para roles corporativos).
\end{itemize}




Cada skill a abordar según necesidad detectada por entrevistas y ofertas: cursos online, tutoriales, prácticas rápidas.




\section{Conclusión y pasos siguientes}




Alexander debe ejecutar este plan de manera disciplinada. En resumen: 




\begin{itemize}
\item \textbf{Posicionamiento:} Presentarse como Technical Artist Unity/WebGL, poniendo énfasis en proyectos concretos (TwinSight X500).  
\item \textbf{Outreach:} Utilizar LinkedIn proactivamente, impulsar su visibilidad en comunidades especializadas, y construir una red de contactos.  
\item \textbf{Herramientas:} Establecer formalidades contractuales (RUT, facturación), y familiarizarse con EOR/plataformas (Deel, Payoneer) para agilizar la contratación.  
\item \textbf{Metas cortas:} Dentro de 3 meses debe tener entrevistas con empresas internacionales, preparando bases para una oferta remota.  
\item \textbf{Metas largas:} Consolidar un empleo que le pague >USD 3k/mes al año, y planificar residencia europea (matrimonio) como respaldo. Seguir mejorando portfolio con nuevos demos y aprendizajes técnicos.  
\end{itemize}




Finalmente, revisar periódicamente el progreso: ajustar roles objetivo según respuestas del mercado, optimizar CV/LinkedIn tras feedback, y mantener actitud analítica (“ingenieril”) en cada paso. Con una estrategia fundamentada en datos actuales y centrada en sus fortalezas técnicas, Alexander puede maximizar sus oportunidades de un empleo internacional bien remunerado en 3–6 meses.



